\documentclass[14pt,a4paper]{extarticle}

\usepackage{cmap}
\usepackage[T2A]{fontenc}
\usepackage[utf8]{inputenc}
\usepackage[russian]{babel}
\usepackage{newtxtext,newtxmath}

\usepackage{amsmath}
\usepackage{geometry}
\usepackage{setspace}
\usepackage{indentfirst}
\usepackage{graphicx}
\usepackage{array}
\usepackage{caption}
\usepackage{float}
\usepackage{booktabs}
\usepackage{siunitx}
\sisetup{output-decimal-marker={,}, separate-uncertainty, per-mode=symbol}
\captionsetup[table]{labelsep=endash, justification=raggedright, singlelinecheck=false}
\captionsetup[figure]{labelsep=endash}
\usepackage{pgfplots}
\pgfplotsset{compat=1.18}

\geometry{left=30mm, right=10mm, top=20mm, bottom=20mm}
\onehalfspacing
\setlength{\parindent}{1.25cm}
\pagestyle{plain}

\begin{document}

\begin{titlepage}
\begin{center}
\onehalfspacing
Белорусский государственный университет\\
Факультет радиофизики и компьютерных технологий

\vspace{4cm}

\textbf{ОТЧЁТ}\\
по лабораторной работе №3\\[0.3cm]
\textbf{«Интерференция света в тонких плёнках»}

\vspace{4cm}

\begin{flushright}
\begin{tabular}{@{}p{7.5cm}@{}}
\raggedright
Выполнил:\\
студент 2 курса 8 группы\\
Нацевский М. А.\\
Преподаватель:\\
Трухан С. В.\\
\end{tabular}
\end{flushright}

\vfill
Минск, 2026
\end{center}
\end{titlepage}
\setcounter{page}{2}

\section*{Цель работы}

Изучить интерференцию света в тонких плёнках на примере полос равной толщины. Определить длины волн отдельных линий излучения ртутной лампы и взаимную деформацию линзы и пластинки.

\section*{Краткая теория}

Кольца Ньютона --- полосы равной толщины, которые возникают в воздушном зазоре между плоской стеклянной пластинкой и выпуклой поверхностью линзы большого радиуса кривизны $R$ (рис.~\ref{fig:rings}). Световые волны, отражённые от верхней и нижней границ зазора, когерентны и интерферируют между собой.

\begin{figure}[htbp]
\centering
\begin{tikzpicture}[scale=0.8]
  \fill[blue!8] (-2.5,0.67) arc[start angle=240, end angle=300, radius=5] -- (2.5,1.7) -- (-2.5,1.7) -- cycle;
  \draw[thick] (-2.5,0.67) arc[start angle=240, end angle=300, radius=5] -- (2.5,1.7) -- (-2.5,1.7) -- cycle;
  \fill[gray!20] (-3.4,-0.6) rectangle (3.4,0);
  \draw[thick] (-3.4,-0.6) rectangle (3.4,0);
  \draw[dashed] (0,-0.6) -- (0,5.2);
  \fill (0,5) circle (1.5pt);
  \node[right] at (0,5.1) {$C$};
  \draw[->, >=stealth] (0,5) -- node[pos=0.3, right] {$R$} (1.8,0.335);
  \draw[<->, >=stealth] (0,-0.95) -- node[below] {$r_m$} (1.8,-0.95);
  \draw[densely dotted] (1.8,0.335) -- (1.8,-0.95);
  \draw[<->, >=stealth] (3.0,0) -- (3.0,0.335);
  \draw[densely dotted] (1.8,0.335) -- (3.0,0.335);
  \node[right] at (3.0,0.45) {$h_m$};
  \draw[red, thick, ->, >=stealth] (-1.6,3.6) -- (-1.6,0.26);
\end{tikzpicture}
\caption{Схема образования колец Ньютона}
\label{fig:rings}
\end{figure}

При нормальном падении оптическая разность хода в месте зазора толщиной $h_m$ равна $2h_m + \lambda/2$, где слагаемое $\lambda/2$ учитывает изменение фазы при отражении от пластинки. Условие тёмного кольца номера $m$:
\begin{equation}
  2h_m + \frac{\lambda}{2} = (2m + 1)\frac{\lambda}{2}.
  \label{eq:dark}
\end{equation}

С учётом взаимной деформации $\Delta h$ линзы и пластинки из геометрии рис.~\ref{fig:rings} следует $2h_m = r_m^2/R - 2\Delta h$. Для диаметра тёмного кольца получаем линейную зависимость от его номера:
\begin{equation}
  d_m^2 = 4R\lambda m + 8R\Delta h = \beta_1 m + \beta_0.
  \label{eq:d2}
\end{equation}

Если диаметры выражены в делениях шкалы окулярного микрометра с ценой деления $a$, то
\begin{equation}
  \lambda = \frac{a^2 \beta_1}{4R}, \qquad \Delta h = \frac{a^2 \beta_0}{8R}.
  \label{eq:lam}
\end{equation}

\section*{Экспериментальная часть}

Излучение ртутной лампы проходило через монохроматор УМ-2 и направлялось в опак-иллюминатор микроскопа, в поле зрения которого наблюдались кольца Ньютона. С помощью окулярного микрометра измерялись левая $x_{\text{л}}$ и правая $x_{\text{п}}$ координаты тёмных колец, диаметр кольца $d_m = x_{\text{л}} - x_{\text{п}}$. Цена деления шкалы $a = 0{,}272$~мм/дел., радиус кривизны линзы $R = 86{,}9$~мм.

Результаты приведены в табл.~\ref{tab:rings}. Для зелёного и жёлтого света измерено по 15 колец, для фиолетово-синего --- 10 колец.

\begin{table}[htbp]
\caption{Координаты тёмных колец (дел.) и квадраты их диаметров (дел.$^2$)}
\label{tab:rings}
\centering
\footnotesize
\setlength{\tabcolsep}{3.5pt}
\begin{tabular}{cccccccccc}
\toprule
& \multicolumn{3}{c}{Зелёный} & \multicolumn{3}{c}{Жёлтый} & \multicolumn{3}{c}{Фиолетово-синий} \\
\cmidrule(lr){2-4} \cmidrule(lr){5-7} \cmidrule(lr){8-10}
$m$ & $x_{\text{л}}$ & $x_{\text{п}}$ & $d_m^2$ & $x_{\text{л}}$ & $x_{\text{п}}$ & $d_m^2$ & $x_{\text{л}}$ & $x_{\text{п}}$ & $d_m^2$ \\
\midrule
1 & 5,34 & 2,74 & 6,76 & 5,37 & 2,71 & 7,08 & 5,85 & 1,77 & 16,65$^{*}$ \\
2 & 5,57 & 2,51 & 9,36 & 5,58 & 2,48 & 9,61 & 6,02 & 1,63 & 19,27 \\
3 & 5,73 & 2,33 & 11,56 & 5,77 & 2,27 & 12,25 & 6,13 & 1,48 & 21,62 \\
4 & 5,91 & 2,16 & 14,06 & 5,94 & 2,11 & 14,67 & 6,22 & 1,40 & 23,23 \\
5 & 6,10 & 2,02 & 16,65 & 6,13 & 1,95 & 17,47 & 6,30 & 1,29 & 25,10 \\
6 & 6,21 & 1,88 & 18,75 & 6,25 & 1,82 & 19,62 & 6,39 & 1,18 & 27,14 \\
7 & 6,33 & 1,75 & 20,98 & 6,38 & 1,68 & 22,09 & 6,48 & 1,09 & 29,05 \\
8 & 6,46 & 1,62 & 23,43 & 6,51 & 1,54 & 24,70 & 6,55 & 1,00 & 30,80 \\
9 & 6,58 & 1,51 & 25,70 & 6,62 & 1,42 & 27,04 & 6,65 & 0,93 & 32,72 \\
10 & 6,70 & 1,39 & 28,20 & 6,73 & 1,31 & 29,38 & 6,74 & 0,88 & 34,34 \\
11 & 6,79 & 1,28 & 30,36 & 6,88 & 1,19 & 32,38 & & & \\
12 & 6,91 & 1,17 & 32,95 & 6,95 & 1,10 & 34,22 & & & \\
13 & 6,97 & 1,09 & 34,57 & 7,05 & 1,00 & 36,60 & & & \\
14 & 7,08 & 0,99 & 37,09 & 7,15 & 0,93 & 38,69 & & & \\
15 & 7,17 & 0,89 & 39,44 & 7,24 & 0,83 & 41,09 & & & \\
\bottomrule
\multicolumn{10}{l}{$^{*}$ Промах, исключён из расчёта.}
\end{tabular}
\end{table}

\section*{Обработка экспериментальных данных}

Зависимости $d_m^2 = \beta_0 + \beta_1 m$ (рис.~\ref{fig:d2}) аппроксимированы методом наименьших квадратов:
\[
  \beta_1 = \frac{N\sum m_i d_{m_i}^2 - \sum m_i \sum d_{m_i}^2}{N\sum m_i^2 - \left(\sum m_i\right)^2},
  \qquad
  \beta_0 = \frac{1}{N}\left(\sum d_{m_i}^2 - \beta_1 \sum m_i\right).
\]

Точки проверялись по стьюдентизированным остаткам с поправкой Бонферрони. Для фиолетово-синего света кольцо №1 отклоняется от прямой значимо (остаток $4{,}85$, вероятность $0{,}019 < 0{,}05$), поэтому оно исключено из расчёта. Для зелёного и жёлтого света промахов не обнаружено.

\begin{figure}[htbp]
\centering
\begin{tikzpicture}
\begin{axis}[
  width=0.9\textwidth, height=9cm,
  xlabel={$m$}, ylabel={$d_m^2$, дел.$^2$},
  grid=major, xmin=0, xmax=16, ymin=0, ymax=45,
  xtick distance=2,
  legend pos=south east, legend cell align=left,
  legend style={font=\small}
]
\addplot[only marks, mark=*, mark size=1.8pt, green!55!black] coordinates {
(1,6.76) (2,9.36) (3,11.56) (4,14.06) (5,16.65) (6,18.75) (7,20.98) (8,23.43)
(9,25.70) (10,28.20) (11,30.36) (12,32.95) (13,34.57) (14,37.09) (15,39.44)
};
\addplot[only marks, mark=square*, mark size=1.6pt, orange!90!black] coordinates {
(1,7.08) (2,9.61) (3,12.25) (4,14.67) (5,17.47) (6,19.62) (7,22.09) (8,24.70)
(9,27.04) (10,29.38) (11,32.38) (12,34.22) (13,36.60) (14,38.69) (15,41.09)
};
\addplot[only marks, mark=triangle*, mark size=2.2pt, violet] coordinates {
(2,19.27) (3,21.62) (4,23.23) (5,25.10) (6,27.14) (7,29.05) (8,30.80) (9,32.72) (10,34.34)
};
\addplot[only marks, mark=o, mark size=3pt, thick, black] coordinates {(1,16.65)};
\addplot[green!55!black, thick, domain=0:15.5, forget plot] {4.739 + 2.3231*x};
\addplot[orange!90!black, thick, domain=0:15.5, forget plot] {4.982 + 2.4346*x};
\addplot[violet, thick, domain=0:10.5, forget plot] {15.767 + 1.8775*x};
\legend{Зелёный, Жёлтый, Фиолетово-синий, Промах}
\end{axis}
\end{tikzpicture}
\caption{Зависимость квадрата диаметра тёмного кольца от его номера}
\label{fig:d2}
\end{figure}

Длины волн и деформация рассчитаны по формулам~\eqref{eq:lam}. Погрешности соответствуют доверительной вероятности $0{,}95$ и учитывают только случайный разброс точек. Результаты приведены в табл.~\ref{tab:res}. Для жёлтого света сравнение выполнено со средним значением линий дублета ртути 5769,60~\AA{} и 5790,66~\AA.

\begin{table}[H]
\caption{Результаты обработки}
\label{tab:res}
\centering
\small
\setlength{\tabcolsep}{4pt}
\renewcommand{\arraystretch}{1.2}
\begin{tabular}{lccccc}
\toprule
Свет & $\beta_1$, дел.$^2$ & $\beta_0$, дел.$^2$ & $\lambda$, \AA & $\lambda_{\text{табл}}$, \AA & $\Delta h$, нм \\
\midrule
Зелёный & $2{,}323 \pm 0{,}027$ & $4{,}74 \pm 0{,}25$ & $4945 \pm 58$ & 5460,72 & $504 \pm 26$ \\
Жёлтый & $2{,}435 \pm 0{,}037$ & $4{,}98 \pm 0{,}34$ & $5182 \pm 78$ & 5780,13 & $530 \pm 36$ \\
Фиолетово-синий & $1{,}878 \pm 0{,}051$ & --- & $3996 \pm 108$ & 4358,34 & --- \\
\bottomrule
\end{tabular}
\end{table}

Для фиолетово-синего света отрезок на оси ординат не использовался. При нумерации колец с единицы он равен $(15{,}77 \pm 0{,}33)$~дел.$^2$, что примерно втрое больше, чем для других цветов. Такое значение соответствует сдвигу нумерации примерно на 6 колец, то есть измерения, по-видимому, начаты не с первого кольца. На наклон прямой и на длину волны этот сдвиг не влияет.

Средневзвешенное значение взаимной деформации по зелёному и жёлтому свету:
\[
  \Delta h = (513 \pm 21)~\text{нм}.
\]

Полученные длины волн меньше табличных на 9,5~\% для зелёного, 10,3~\% для жёлтого и 8,3~\% для фиолетово-синего света. Это отклонение превышает случайную погрешность и почти одинаково для всех цветов, что указывает на систематическую ошибку масштаба. Отношение наклонов $\beta_1^{\text{ф}}/\beta_1^{\text{з}} = 0{,}808$, не зависящее от $R$ и $a$, близко к отношению табличных длин волн $0{,}798$. Значение $R$ указано на установке, а цена деления $a$ зависит от объектива микроскопа, поэтому вероятной причиной является несоответствие принятого значения $a$ используемому объективу. Согласие с табличными значениями достигалось бы при $a \approx 0{,}286$~мм/дел.

\section*{Вывод}

Изучена интерференция света в тонкой воздушной прослойке на примере колец Ньютона. Подтверждена линейная зависимость квадрата диаметра тёмного кольца от его номера. Получены длины волн $\lambda_{\text{з}} = (4945 \pm 58)$~\AA, $\lambda_{\text{ж}} = (5182 \pm 78)$~\AA{} и $\lambda_{\text{ф}} = (3996 \pm 108)$~\AA, которые на 8--10~\% меньше табличных из-за систематической ошибки масштаба. Взаимная деформация линзы и пластинки составила $\Delta h = (513 \pm 21)$~нм.

\end{document}
